\section{Conclusion}

We presented a systematic evaluation of long-horizon auto research agents that goes beyond final scores by examining Solution Framing, Execution, Feedback Control, idea-level novelty, experience reuse, and harness effects.
The results place current systems at a stage of partial research-loop automation.
Agents can identify practical approaches, implement them, and sometimes reach competitive solutions, but their strongest behavior is not reproduced consistently, and genuine innovation remains rare.
Process bottlenecks vary across models and workloads, experience transfer remains unstable, and harnesses mainly improve the reliable realization of existing capability.
These distinctions map observed limitations to targeted training, inference-time selection, selective memory, harness design, or stronger tasks and verifiers.
The resulting evaluation provides a concrete basis for improving auto research agents and tracking progress toward more reliable and self-improving research systems.
